% Created 2026-01-13 Tue 22:33
% Intended LaTeX compiler: lualatex
\documentclass[11pt]{article}
\usepackage{minted}
\usepackage{amsmath}
\usepackage{fontspec}
\usepackage{graphicx}
\usepackage{longtable}
\usepackage{wrapfig}
\usepackage{rotating}
\usepackage[normalem]{ulem}
\usepackage{capt-of}
\usepackage{hyperref}
\usepackage{minted}
\usepackage{amssymb}
\usepackage{xcolor}
\usepackage{newunicodechar}
\usepackage{newunicodechar}
\newunicodechar{₁}{\ensuremath{_1}}
\newunicodechar{₂}{\ensuremath{_2}}
\newunicodechar{₃}{\ensuremath{_3}}
\newunicodechar{₄}{\ensuremath{_4}}
\newunicodechar{₅}{\ensuremath{_5}}
\newunicodechar{₆}{\ensuremath{_6}}
\newunicodechar{₇}{\ensuremath{_7}}
\newunicodechar{₈}{\ensuremath{_8}}
\newunicodechar{₉}{\ensuremath{_9}}
\newunicodechar{₀}{\ensuremath{_0}}
\newunicodechar{ₐ}{\ensuremath{_a}}
\newunicodechar{ₑ}{\ensuremath{_e}}
\newunicodechar{ₒ}{\ensuremath{_o}}
\newunicodechar{ₓ}{\ensuremath{_x}}
\newunicodechar{ₙ}{\ensuremath{_n}}
\newunicodechar{ₖ}{\ensuremath{_k}}
\newunicodechar{ₘ}{\ensuremath{_m}}
\newunicodechar{≡}{\ensuremath{\equiv}}
\newunicodechar{…}{\ensuremath{\ldots}}
\newunicodechar{≈}{\ensuremath{\mathrel{\approx}}}
\newunicodechar{↓}{\ensuremath{\mathrel{\downarrow}}}
\newunicodechar{✅}{\textcolor{green!50!black}{\ensuremath{\checkmark}}}
\newunicodechar{❌}{\textcolor{red!70!black}{\ensuremath{\times}}}
\newunicodechar{∣}{\ensuremath{\mid}}
\newunicodechar{∈}{\ensuremath{\in}}
\newunicodechar{→}{\ensuremath{\to}}
\newunicodechar{⇒}{\ensuremath{\Rightarrow}}
\newunicodechar{↦}{\ensuremath{\mapsto}}
\newunicodechar{≥}{\ensuremath{\ge}}
\newunicodechar{≤}{\ensuremath{\le}}
\newunicodechar{≡}{\ensuremath{\equiv}}
\newunicodechar{⊂}{\ensuremath{\subset}}
\newunicodechar{⊃}{\ensuremath{\supset}}
\newunicodechar{⊆}{\ensuremath{\subseteq}}
\newunicodechar{⊇}{\ensuremath{\supseteq}}
\newunicodechar{⟶}{\ensuremath{\longrightarrow}}
\newunicodechar{⟵}{\ensuremath{\longleftarrow}}
\newunicodechar{⋯}{\ensuremath{\cdots}}
\newunicodechar{…}{\ensuremath{\ldots}}
\newunicodechar{∪}{\ensuremath{\cup}}
\newunicodechar{∩}{\ensuremath{\cap}}
\newunicodechar{∧}{\ensuremath{\wedge}}
\newunicodechar{∨}{\ensuremath{\vee}}
\newunicodechar{¬}{\ensuremath{\neg}}
\newunicodechar{∀}{\ensuremath{\forall}}
\newunicodechar{∃}{\ensuremath{\exists}}
\newunicodechar{⊢}{\ensuremath{\vdash}}
\newunicodechar{⊨}{\ensuremath{\vDash}}
\newunicodechar{⊥}{\ensuremath{\bot}}
\newunicodechar{⊤}{\ensuremath{\top}}
\newunicodechar{∅}{\ensuremath{\emptyset}}
\newunicodechar{↔}{\ensuremath{\leftrightarrow}}
\newunicodechar{⇔}{\ensuremath{\Leftrightarrow}}
\usepackage[
paper=letterpaper,
left=1in,
right=1in,
top=1in,
bottom=1in
]{geometry}
\usepackage{fontspec}
\setmainfont{Avenir Next}[
UprightFont = *,
ItalicFont = * Italic,
BoldFont = * Demi Bold,
BoldItalicFont = * Demi Bold Italic
]
\author{Michael Hohn}
\date{\today}
\title{}
\hypersetup{
 pdfauthor={Michael Hohn},
 pdftitle={},
 pdfkeywords={},
 pdfsubject={},
 pdfcreator={Emacs 30.2 (Org mode 9.7.11)}, 
 pdflang={English}}
\begin{document}

\tableofcontents

\section{MRVA documentation: system definitions}
\label{sec:org8d205f4}
This document defines the MRVA components, non-components, and their relation
to development, deployment, and documentation.

In summary:

\begin{enumerate}
\item MRVA core applies existing CodeQL queries to existing CodeQL databases
using the CodeQL command-line tools, and produces SARIF results.
It performs no query authoring, result interpretation, or presentation.

\item The only valid interface to the MRVA core consists of exactly three
requests, issued via \texttt{gh-mrva}:
\begin{itemize}
\item submit
\item status
\item get-results
\end{itemize}

No other execution interfaces are supported.

\item User interfaces are explicitly outside the core and exist solely for
convenience.
For demonstration of the core interface (developers) and its operation
(users), simple Tk-based UIs are provided.
These interfaces are non-authoritative and replaceable.

\item CodeQL database provisioning is required by the MRVA core but is not part
of the core itself.
The MRVA core depends on the existence of a database provider role, not on
any specific implementation.

The database provider interface consists of a single operation:
\texttt{list-dbs}.

A reference implementation of this interface is provided by the \texttt{hepc}
service.
This implementation targets open-source CodeQL databases and exists to:
\begin{itemize}
\item demonstrate the provider interface
\item enable development and testing
\item provide a complete, runnable system
\end{itemize}

The \texttt{hepc} service requires substantial engineering effort but remains
strictly outside the MRVA core.

In deployed environments, the reference implementation is expected to be
replaced by site-specific database providers (for example, internal or
proprietary repositories) without changes to the MRVA core.
\end{enumerate}
\subsection{MRVA Core definition}
\label{sec:org67b83e3}
MRVA core is the minimal, sufficient system that turns
\begin{verbatim}
(CodeQL databases × queries) → SARIF results
\end{verbatim}

Nothing else.

\begin{itemize}
\item The core
\begin{itemize}
\item Accepts references to
\begin{itemize}
\item existing CodeQL databases
\item existing CodeQL queries (packs or individual queries)
\end{itemize}
\item Executes CodeQL deterministically
\item Produces SARIF output
\item Exposes execution state and results
\end{itemize}

\item The core does not:
\begin{itemize}
\item edit queries
\item render results
\item interpret findings
\item provide rich UIs
\item manage source code checkouts
\item manage DB construction
\end{itemize}

\item MRVA core is an execution engine, not an IDE, not an analyst tool.
\end{itemize}
\subsection{MRVA Core Interface, Authoritative API}
\label{sec:orgb53e0dc}
\begin{itemize}
\item Interface Name: \texttt{gh-mrva}
This is the only supported interface to the core.

\item Interface Style
\begin{itemize}
\item Shell-driven
\item Scriptable
\item Stateless per invocation
\end{itemize}

\item Outside the core, this interface may be used:
\begin{itemize}
\item directly via the shell
\item by user interfaces
\item by automated systems (including AI agents)
\end{itemize}

\item Command contract:
Exactly three requests are supported.
\begin{itemize}
\item submit
\begin{itemize}
\item Inputs:
\begin{itemize}
\item DB identifiers
\item Query identifiers
\item Optional execution parameters
\end{itemize}
\item Output:
\begin{itemize}
\item opaque job ID
\end{itemize}
\end{itemize}

\item status
\begin{itemize}
\item Input:
\begin{itemize}
\item job ID
\end{itemize}
\item Output:
\begin{itemize}
\item state ∈ \{queued, running, finished, failed\}
\item minimal progress metadata
\end{itemize}
\end{itemize}

\item get-results
\begin{itemize}
\item Input:
\begin{itemize}
\item job ID
\end{itemize}
\item Output:
\begin{itemize}
\item SARIF results
\item CodeQL databases (optional, if explicitly requested)
\end{itemize}
\end{itemize}
\end{itemize}
\end{itemize}
\subsection{MRVA Convenience UI for gh-mrva (Optional, Replaceable)}
\label{sec:orgf19b4af}
\begin{itemize}
\item Purpose:  Reduce friction for humans.

\item Implementations
\begin{itemize}
\item Basic UI (Python / Tk)
\begin{itemize}
\item thin wrapper over gh-mrva
\item submit jobs
\item poll status
\item fetch SARIF
\item minimal result listing

\item Properties:
\begin{itemize}
\item disposable
\item non-authoritative
\item allowed to be ugly
\end{itemize}
\end{itemize}

\item Fancy UI (VS Code extension)
\begin{itemize}
\item strictly a gh-mrva client
\item no embedded execution logic
\item no hidden or private APIs
\end{itemize}
\end{itemize}

\item Key rule:
If a VS Code extension requires more privileges or APIs than shell access
to gh-mrva, the design is incorrect.
\end{itemize}
\subsection{MRVA Core Result Handler (Post-Processing Layer)}
\label{sec:org35c5525}
\begin{itemize}
\item Purpose:
Transform raw SARIF into a structured, queryable dataset.

\item Responsibilities:
\begin{itemize}
\item ingest SARIF
\item normalize fields
\item load data into DuckDB
\item preserve references to:
\begin{itemize}
\item repository
\item database
\item query
\item execution metadata
\end{itemize}
\end{itemize}

\item Output:
A read-only analytical store.

\item Dependency rules:
\begin{itemize}
\item MRVA core does not depend on DuckDB
\item DuckDB ingestion depends on MRVA output
\end{itemize}
\end{itemize}
\subsection{Non-core: Query Editing}
\label{sec:org10625ca}
\begin{itemize}
\item Scope: Query editing is not MRVA.

\item Location
\begin{itemize}
\item VS Code
\item Plain filesystem
\item Existing CodeQL tooling
\end{itemize}

\item Relationship to MRVA:
\begin{itemize}
\item MRVA:
\begin{itemize}
\item consumes queries
\item does not assist with authoring
\item does not validate intent
\end{itemize}
\end{itemize}

\item This separation avoids:
\begin{itemize}
\item IDE entanglement
\item schema drift
\item toolchain coupling
\end{itemize}
\end{itemize}
\subsection{Development}
\label{sec:org0f0dabe}
\begin{itemize}
\item Definition: Development is everything that exists to change MRVA.

\item Scope Includes
\begin{itemize}
\item source code
\item build scripts
\item test scaffolding
\item experimental tools
\item refactors
\item prototype UIs
\item temporary glue
\end{itemize}

\item Properties
\begin{itemize}
\item Mutable
\item Incomplete
\item Allowed to be ugly
\item Allowed to be slow
\item Allowed to break
\end{itemize}

\item Key Rule: Nothing in development is relied upon at runtime unless explicitly
promoted to deployment.
\end{itemize}
\subsection{Deployment}
\label{sec:org8f8aecc}
\begin{itemize}
\item Definition: Deployment is the sealed, operated system.
This is what users run.

\item Scope includes:
\begin{itemize}
\item MRVA core binaries or containers
\item gh-mrva
\item fixed execution environments (Linux)
\item deterministic configuration
\end{itemize}

\item Properties:
\begin{itemize}
\item immutable (by policy)
\item reproducible
\item predictable
\item auditable
\end{itemize}

\item Key rule:
Deployment artifacts must be buildable from development using a fixed,
documented procedure.
\end{itemize}
\subsection{Documentation}
\label{sec:orgc6b6f12}
Documentation is split along the following axes:
\begin{itemize}
\item system: definitions, structure
\item development: setup, use
\item deployment: setup, use
\end{itemize}
\end{document}
